\chapter{Introduction}
\label{ch:intro}

This chapter states the clinical problem that MedLift-3D solves and why it matters. It then lists the objectives as testable hypotheses, summarises the method and its outcomes, and outlines the rest of the report.

\section{Project Overview}

Lung cancer is the leading cause of cancer death worldwide, with about 2.5 million new cases (12.4\% of all cancers) and 1.8 million deaths (18.7\% of cancer deaths) in 2022~\cite{bray2024globocan}. Survival depends strongly on the stage at diagnosis: five-year survival is 70--90\% for stage~I non-small cell lung cancer but only 15--19\% for stage~IV~\cite{knight2017early}. Early lung cancer usually appears as a pulmonary nodule, yet most nodules are benign. Radiologists therefore decide by following a nodule over months and measuring how fast it grows, usually as a volume doubling time (VDT)~\cite{revel2004threed,horeweg2014nelson}.

Computed tomography (CT) is the standard tool for this follow-up, but every scan gives the patient a radiation dose of several millisieverts and costs money and scanner time~\cite{brenner2007ct}. Low-dose CT lowers the dose, yet repeated follow-up still adds up: about four radiation-induced cancers per \num{10000} screened people, most of it from the follow-up scans~\cite{rampinelli2017}. In low-resource settings such as rural Bangladesh, CT is often not available at all.

Digital chest tomosynthesis (DTS) acquires 15--74 low-dose X-ray views during a short tube sweep of about \ang{30}, at roughly \SI{0.1}{mSv}, on equipment that costs far less than CT~\cite{dobbins2003dts}. It finds nodules about three times more often than chest radiography~\cite{vikgren2008}. However, a \ang{30} sweep leaves a \emph{missing wedge} of directions that no X-ray passes through (\cref{fig:dts}), so DTS cannot resolve depth and every 3D size measured from it is partly a guess.

\paragraph{Problem statement.} Given a low-dose chest DTS acquisition and one click on a nodule, produce its 3D mass, volume and growth rate together with an uncertainty that is honest, so that a clinician can follow nodules without repeated CT and knows when a measurement cannot be trusted.

\begin{figure}[htbp]
  \centering
  \resizebox{\linewidth}{!}{% DTS acquisition geometry (a) and the missing wedge in Fourier space (b)
\begin{tikzpicture}[font=\footnotesize, >=Stealth]
  % ---------------- (a) geometry
  \begin{scope}
    % detector
    \fill[edge!25] (-3.6,-0.18) rectangle (3.6,0);
    \draw[edge] (-3.6,-0.18) rectangle (3.6,0);
    \node[below=1pt] at (0,-0.18) {flat-panel detector (43 cm, 0.8 mm pixels)};
    % ray fans (drawn first, behind the patient)
    \foreach \a in {-15, 0, 15} {
      \fill[uiublue, opacity=0.07] ({6.2*tan(\a)},6.2) -- (-3.6,0) -- (3.6,0) -- cycle;
    }
    % patient cross-section
    \draw[edge, fill=physics] (0,1.25) ellipse (2.7 and 1.05);
    \draw[edge!60, fill=white] (-1.25,1.3) ellipse (0.95 and 0.68);
    \draw[edge!60, fill=white] (1.25,1.3) ellipse (0.95 and 0.68);
    % source track
    \draw[edge, dashed] (-2.4,6.2) -- (2.4,6.2);
    \foreach \a in {-15,-12.857,...,15.1} {
      \fill[edge] ({6.2*tan(\a)},6.2) circle (0.035);
    }
    % rays through the nodule to the detector: the shadow moves with the tube angle
    \foreach \a in {-15, 0, 15} {
      \coordinate (S) at ({6.2*tan(\a)},6.2);
      \draw[accent, thick] (S) -- ($(S)!{6.2/(6.2-1.45)}!(0.95,1.45)$) coordinate (D);
      \fill[accent] (D) circle (0.06);
      \fill[black] (S) ++(-0.14,0) rectangle ++(0.28,0.22);
    }
    \fill[accent] (0.95,1.45) circle (0.09);
    \node[accent, right=1pt, fill=physics, inner sep=1pt] at (1.07,1.62) {nodule};
    \node[above=3pt, align=center] at (0,6.4) {X-ray tube: 15 positions on a linear track};
    % sweep angle, measured at the detector centre
    \draw[edge, densely dotted] (0,0) -- (0,6.2);
    \draw[edge, densely dotted] (0,0) -- ({6.2*tan(15)},6.2);
    \draw[edge, densely dotted] (0,0) -- ({-6.2*tan(15)},6.2);
    \draw[edge, <->] ({4.4*cos(105)},{4.4*sin(105)}) arc[start angle=105, end angle=75, radius=4.4];
    \node[fill=white, inner sep=1pt] at (0,4.65) {$\pm 15^\circ$};
    % SID dimension
    \draw[<->, edge] (3.95,0) -- (3.95,6.2) node[midway, right, align=left] {SID\\1800 mm};
    % axes
    \draw[->] (-3.95,0.3) -- (-3.95,1.6) node[above] {depth $y$};
    \draw[->] (-3.95,0.3) -- (-2.9,0.3) node[right] {$x$};
    \node[font=\small\bfseries] at (-3.7,6.9) {(a)};
  \end{scope}

  % ---------------- (b) Fourier space
  \begin{scope}[xshift=8.3cm, yshift=3.0cm]
    \def\R{2.1}
    \fill[uiublue!25] (0,0) -- (-15:\R) arc[start angle=-15, end angle=15, radius=\R] -- cycle;
    \fill[uiublue!25] (0,0) -- (165:\R) arc[start angle=165, end angle=195, radius=\R] -- cycle;
    \fill[pattern=north east lines, pattern color=red!60] (0,0) -- (15:\R) arc[start angle=15, end angle=165, radius=\R] -- cycle;
    \fill[pattern=north east lines, pattern color=red!60] (0,0) -- (195:\R) arc[start angle=195, end angle=345, radius=\R] -- cycle;
    \draw[edge] (0,0) circle (\R);
    \draw[->] (-2.5,0) -- (2.6,0) node[right] {$k_x$};
    \draw[->] (0,-2.5) -- (0,2.6) node[above] {$k_y$ (depth)};
    \node[fill=white, inner sep=1pt, align=center] at (0,1.35) {missing wedge\\(not measured)};
    \node[align=center, text=uiublue!70!black] at (1.25,-0.38) {measured};
    \draw[edge] (1.3,0) arc[start angle=0, end angle=15, radius=1.3];
    \node at (1.62,0.2) {$15^\circ$};
    \node[font=\small\bfseries] at (-2.6,3.9) {(b)};
  \end{scope}
\end{tikzpicture}
}
  \caption[DTS acquisition geometry and the missing wedge]{Chest DTS acquisition. (a) The X-ray tube moves along a linear track over $\pm15^\circ$ and the nodule's shadow shifts on the detector with the tube angle; this shift is the only depth information. (b) In Fourier space each view measures one plane, so a $\pm15^\circ$ sweep measures only a double wedge around the in-plane axis and leaves the depth frequencies unmeasured.}
  \label{fig:dts}
\end{figure}

\section{Motivation}

This project is motivated by four observations: nodules need to be followed over time, the size of nodules is uncertain when measured in 2D, follow-up using CT is costly because of radiation, access, and existing reconstruction methods are not able to produce trustworthy numbers from DTS. Our approach is guided by a 5th, computational observation.

\subsection{Clinical Importance of Nodule Follow-up}
The 5-year survival rates are 64\% for localised lung cancer, and 9\% for distant disease~\cite{who2023lung,nlst2011}. The only distinguishing feature between malignant and benign nodules is growth and the guidelines are based on ranges of volume and doubling times: $<100$, 100--300 and $\geq$300 mm$^3$; $<400$, 400--600 and $>600$ days~\cite{horeweg2014nelson,oudkerk2017eups}.

\subsection{Limits of 2D Measurement}
 Diameters measured on single slices assume a sphere. Volumes estimated from 2D measurements differ from 3D volumetry by 47--85\% on average~\cite{demargerie2018}, 2D measurements of small nodules are unreliable~\cite{revel2004twod}, and a 1~mm change in diameter is a 33\% change in volume~\cite{lee2025volumetry}. Measurement error larger than the true growth delays treatment.

\subsection{Radiation and Access}
 CT follow-up means repeated dose and requires well-equipped centres~\cite{mascalchi2017,brenner2007ct}. In Bangladesh, CT scanners are concentrated in divisional cities and large hospitals, so many patients travel long distances and pay for every follow-up scan, while digital radiography rooms, which can be upgraded to tomosynthesis, are far more widespread. DTS is cheaper, faster and much lower in dose than CT, but today it offers no trustworthy 3D measurement, so a nodule found on DTS is still sent to CT for every follow-up.

\subsection{Why Existing Reconstruction and AI Are Not Enough}
Iterative reconstruction from a \ang{30} sweep results in 3--5 times nodule extension along the depth axis in comparison to in-plane (the median was 3.5 in our measurements of LIDC nodules). The missing wedge is generative models, but a cancer-following model that creates anatomy is hazardous~\cite{bhadra2021,zhang2025pail}. Sparse view methods have been published with image quality (PSNR,SSIM) values not calibrated clinical numbers and none is developed for DTS geometry.

\subsection{Computational Motivation}
Not only what it can show, but we asked what DTS can \emph{measure}. By derivation of projection moments (\cref{sec:method}) it is shown that every DTS view measures the total excess attenuation of a nodule exactly, its depth extent, which volume needs, is about 15 times worse conditioned for a $\pm15^\circ$ sweep. A system for DTS should therefore start with a statement of \emph{mass}, then a fill depth with a learned prior that is consistent with the X-rays, and a statement of uncertainty for each of the numbers. The use of mass growth as a CT biomarker is already proven to detect growth earlier than volume~\cite{dehoop2010}.

\section{Objectives}

The objectives are stated as testable hypotheses (H1--H6); any of them could turn out false, and \cref{ch:results} reports what we found.

\begin{enumerate}[label=\textbf{O\arabic*}]
  \item \textbf{Mass-first measurement (H1).} Show, in theory and on data, that nodule mass from DTS has a smaller error than nodule volume.
  \item \textbf{Learned prior for the missing wedge (H2).} Train a conditional 3D diffusion model that, combined with physics data consistency and 3D Gaussians, reduces volume error compared with physics-only reconstruction.
  \item \textbf{Calibrated uncertainty (H3, H5).} Give every mass and volume a conformal 90\% prediction interval that reaches its nominal coverage on held-out patients, and give calibrated probabilities of the clinical size categories with coverage in every size group.
  \item \textbf{Growth (H4).} Measure growth on virtual follow-up scans and test whether mass doubling time (MDT) classifies growth more accurately than VDT.
  \item \textbf{Explainability (H6).} Show for each nodule which part of the 3D result was measured by the X-rays (seen) and which was supplied by the prior (unseen), and test whether the prior's unseen content recovers the truth rather than inventing it.
  \item \textbf{Usable system.} Deliver a click-to-3D application with a learned detector that suggests where to click, running on free or low-cost cloud hardware.
\end{enumerate}

\section{Methodology}

MedLift-3D is evaluated as a virtual clinical trial on real patient CT, because no public dataset has chest DTS together with CT ground truth. \Cref{fig:pipeline} in \cref{ch:design} shows the full pipeline.
\begin{enumerate}
  \item \textbf{Data:} 583 LIDC-IDRI patients with \num{1144} nodules of 4--30~mm marked by at least three of four radiologists~\cite{armato2011lidc,hancock2016pylidc}, split by patient 60/20/20 into training, calibration and test sets.
  \item \textbf{DTS simulation:} 15 low-dose views over $\pm15^\circ$ with Poisson noise, reconstructed by SIRT~\cite{gilbert1972sirt}.
  \item \textbf{Click and prism:} the clinician, helped by a learned detector, selects a nodule; a 64~mm column through the whole chest (the \emph{prism}) is reconstructed around it.
  \item \textbf{Diffusion prior:} a 17.4-million-parameter conditional 3D U-Net diffusion model~\cite{ho2020ddpm,song2021ddim,salimans2022v} generates eight CT-like hypotheses of the prism, each kept consistent with the X-rays by SIRT steps inside the sampler.
  \item \textbf{3D Gaussians:} each hypothesis is fitted with anisotropic 3D Gaussians~\cite{kerbl2023gs,zha2024r2}, which give the nodule's mass in closed form.
  \item \textbf{Uncertainty:} the spread of the eight hypotheses is calibrated into conformal intervals and predictive distributions on the calibration patients~\cite{vovk2022alrw,vovk2019cps,bostrom2022crepes,angelopoulos2023}.
  \item \textbf{Explanation:} each result is split into seen and unseen parts by a conjugate-gradient null-space decomposition~\cite{bhadra2021,wang2023ddnm}, and the detector is explained with input$\times$gradient saliency~\cite{shrikumar2017}.
\end{enumerate}

\section{Project Outcome}

The project generates the following three sets of outputs: a technical system that is shown to operate; a system usable by a clinician; and a research contribution.

\subsection{Technical Outcome}
An automated, step-by-step, scripted process from data download to final tables, figures and HTML report, which has been verified and passed. It includes a differentiatable DTS simulator and projector, whole-chest and prism reconstruction, a trained conditional diffusion prior and nodule detector, Gaussian measurement, conformal calibration and an explainability module, with 69 unit tests documented and with free use of cloud GPUs via a rental service and Kaggle notebooks.

\subsection{Clinical Outcome}
 A click-to-3D application that returns, in under a minute, the nodule's mass, volume, long-axis diameter and EUPS size category with calibrated probabilities, together with a virtual follow-up for growth. Every number comes with an interval instead of a single value, and a low-confidence flag tells the clinician when the DTS evidence is too weak and a CT is advisable. The tool is meant for following known nodules, the step that today requires repeated CT.

\subsection{Research Contribution}
 (i) A projection-moment argument that DTS measures mass but not depth extent; (ii) a data-consistent diffusion prior inside a DTS prism combined with Gaussian mass; (iii) calibrated, decision-level uncertainty for DTS volumetry; (iv) a truth-referenced seen/unseen test of what the prior adds. \Cref{sec:gap} states exactly what is and is not new.

\section{Organization of the Report}

\Cref{ch:background} provides the background of DTS, diffusion models, 3D Gaussians, conformal prediction and explainability, discusses related applications and research, and identifies the gap. The requirements, context and data flow diagrams, user interface, detailed method, including alternatives examined, project plan and task allocation are presented in \cref{ch:design}. \Cref{ch:results} describes the environment and tools as well as how each requirement was verified and the results of H1-H6. \Cref{ch:standards} discusses standards, issues related to society, health and safety, law, culture, ethics, environment and sustainability, costs and the complex engineering problem mapping. In \cref{ch:conclusion} the work is summarised, limitations are pointed out, future work is suggested and what has been learnt is summarised.
